% Experiment Template
% 1) Copy this file and rename it, e.g. experiments/experiment_04_firewall_rules.tex
% 2) In main.tex, add:
%    \includeexperiment{XX}{Experiment Title}{experiments/your_new_file_name}
% 3) Keep this file content as sections only (no \experiment command here)

\section{Theory}
The Autokey cipher is a substitution cipher that incorporates the message itself into the key. It begins with a seed key (a single character or a word), and the subsequent key characters are taken from the plaintext characters. This prevents the key from repeating, making it more resistant to frequency analysis than the Vigenère cipher. 

The mathematical formula for encryption is:
$$C_i = (P_i + K_i) \mod 26$$

where $K_0$ is the seed key, and $K_i = P_{i-1}$ for $i > 0$. For decryption, the formula is:
$$P_i = (C_i - K_i + 26) \mod 26$$

where the key $K_i$ is derived from previously decrypted plaintext characters.

The seed key is used to shift the first plaintext character, then each subsequent character is shifted by the previous plaintext character's value. 

\begin{table}[H]
  \centering
  \begin{tabularx}{\textwidth}{|*{6}{>{\centering\arraybackslash}X|}}
    \hline
    \textbf{Plaintext} & h & e & l & l & o \\
    \hline
    \textbf{Plaintext Value} & 7 & 4 & 11 & 11 & 14 \\
    \hline
    \textbf{Key} & 3 & 7 & 4 & 11 & 11 \\
    \hline
    \textbf{Ciphertext Value} & 10 & 11 & 15 & 22 & 25 \\
    \hline
    \textbf{Ciphertext} & k & l & p & w & z \\
    \hline
  \end{tabularx}
  \caption{Autokey Cipher Encryption Example}
\end{table}

The first character `h' (7) is shifted by seed key `d' (3): $(7 + 3) \mod 26 = 10$ → `k'. The next character `e' (4) is shifted by the previous plaintext value (7): $(4 + 7) \mod 26 = 11$ → `l'. This process continues with each character shifted by the previous plaintext character's value.

\section{Algorithm}
The Encryption algorithm steps are:
\begin{enumerate}
  \item Input Plaintext and seed key $k$.
  \item Initialize $curr\_k := k$ and index $i := 0$.
  \item While $i < \text{length}(plaintext)$ do:
  \begin{enumerate}
    \item If $char[i]$ is a letter:
    \begin{enumerate}
      \item $P := \text{index of } char[i]$ (0-25).
      \item $char[i] := (P + curr\_k) \mod 26$.
      \item $curr\_k := P$ (Update key with plaintext).
    \end{enumerate}
    \item Increment $i$.
  \end{enumerate}
  \item Return the encrypted text.
\end{enumerate}

The Decryption algorithm steps are:
\begin{enumerate}
  \item Input Ciphertext and seed key $k$.
  \item Initialize $curr\_k := k$ and index $i := 0$.
  \item While $i < \text{length}(ciphertext)$ do:
  \begin{enumerate}
    \item If $char[i]$ is a letter:
    \begin{enumerate}
      \item $C := \text{index of } char[i]$ (0-25).
      \item $P := (C - curr\_k + 26) \mod 26$.
      \item $char[i] := P$.
      \item $curr\_k := P$ (Update key with decrypted plaintext).
    \end{enumerate}
    \item Increment $i$.
  \end{enumerate}
  \item Return the decrypted text.
\end{enumerate}

% \section{Flowchart}
% \begin{figure}[H]
%   \centering
%   \begin{autofitflowchart}
%     \node (start) [terminator] {Start};
%     \node (inputtext) [io, left=of start] {Input Plaintext, Key};
%     \node (initenc) [process, left=of inputtext] {$curr\_k := key$ \\ $i := 0$};
%     \node (checkenc) [decision, below=of initenc] {$i < length$?};
%     \node (checkchar) [decision, right=of checkenc, xshift=2cm] {Is Letter?};
%     \node (shiftenc1) [process, right=of checkchar, xshift=2cm] {$P := char[i]$ \\ $char[i] := (P + curr\_k) \% 26$};
%     \node (shiftenc2) [process, below=of shiftenc1] {$curr\_k := P$};
%     \node (incenc) [process, below=of checkchar] {$i := i + 1$};
%     \node (displayenc) [io, below=of checkenc] {Display Ciphertext};
%     \node (end) [terminator, below=of displayenc] {End};

%     \draw [line] (start) -- (inputtext);
%     \draw [line] (inputtext) -- (initenc);
%     \draw [line] (initenc) -- (checkenc);
%     \draw [line] (checkenc.east) -- node[above]{Yes} (checkchar.west);
%     \draw [line] (checkchar.east) -- node[above]{Yes} (shiftenc1.west);
%     \draw [line] (shiftenc1) -- (shiftenc2);
%     \draw [line] (checkchar.south) -- node[right]{No} (incenc);
%     \draw [line] (shiftenc2.west) -- (incenc.east);
%     \draw [line] (incenc) -- (checkenc);
%     \draw [line] (checkenc) -- node[left]{No} (displayenc);
%     \draw [line] (displayenc) -- (end);
%   \end{autofitflowchart}
%   \caption{Flowchart for Autokey Cipher Encryption}
% \end{figure}

% \begin{figure}[H]
%   \centering
%   \begin{autofitflowchart}
%     \node (start) [terminator] {Start};
%     \node (inputtext) [io, left=of start] {Input Ciphertext, Key};
%     \node (initenc) [process, left=of inputtext] {$curr\_k := key$ \\ $i := 0$};
%     \node (checkenc) [decision, below=of initenc] {$i < length$?};
%     \node (checkchar) [decision, right=of checkenc, xshift=2cm] {Is Letter?};
%     \node (shiftenc1) [process, right=of checkchar, xshift=2cm] {$C := char[i]$ \\ $char[i] := (C - curr\_k + 26) \% 26$};
%     \node (shiftenc2) [process, below=of shiftenc1] {$curr\_k := char[i]$};
%     \node (incenc) [process, below=of checkchar] {$i := i + 1$};
%     \node (displayenc) [io, below=of checkenc] {Display Plaintext};
%     \node (end) [terminator, below=of displayenc] {End};

%     \draw [line] (start) -- (inputtext);
%     \draw [line] (inputtext) -- (initenc);
%     \draw [line] (initenc) -- (checkenc);
%     \draw [line] (checkenc.east) -- node[above]{Yes} (checkchar.west);
%     \draw [line] (checkchar.east) -- node[above]{Yes} (shiftenc1.west);
%     \draw [line] (shiftenc1) -- (shiftenc2);
%     \draw [line] (checkchar.south) -- node[right]{No} (incenc);
%     \draw [line] (shiftenc2.west) -- (incenc.east);
%     \draw [line] (incenc) -- (checkenc);
%     \draw [line] (checkenc) -- node[left]{No} (displayenc);
%     \draw [line] (displayenc) -- (end);
%   \end{autofitflowchart}
%   \caption{Flowchart for Autokey Cipher Decryption}
% \end{figure}

\section{Code}
\begin{minted}{cpp}
#include<bits/stdc++.h>
using namespace std;

class autokey {
  private:
    int key;
  public:
    autokey(char k);
    string encrypt(string s);
    string decrypt(string s);
};

autokey::autokey(char k) {
  while (!isalpha(k)) {
    cout << "Key is not an alphabet. Enter a new key: ";
    cin >> k;
    cin.ignore();
  }
  if (isupper(k)) {
    key = k - 'A';
  } else {
    key = k - 'a';
  }
}

string autokey::encrypt(string s) {
  string ans = s;
  int n = s.length();
  char current_key = key;
  for (int i = 0; i < n; i++) {
    if (s[i] >= 'a' && s[i] <= 'z') {
      ans[i] = ((s[i] - 'a') + current_key) % 26 + 'a';
      current_key = s[i] - 'a';
    } else if (s[i] >= 'A' && s[i] <= 'Z') {
      ans[i] = ((s[i] - 'A') + current_key) % 26 + 'A';
      current_key = s[i] - 'A';
    } else {
      ans[i] = s[i];
    }
  }
  return ans;
}

string autokey::decrypt(string s) {
  string ans = s;
  int n = s.length();
  char current_key = key;
  for (int i = 0; i < n; i++) {
    if (s[i] >= 'a' && s[i] <= 'z') {
      ans[i] = ((s[i] - 'a') - current_key + 26*2) % 26 + 'a';
      current_key = ans[i] - 'a';
    } else if (s[i] >= 'A' && s[i] <= 'Z') {
      ans[i] = ((s[i] - 'A') - current_key + 26*2) % 26 + 'A';
      current_key = ans[i] - 'A';
    } else {
      ans[i] = s[i];
    }
  }
  return ans;
}

int main() {
  cout << "Enter a character key: ";
  char c;
  cin >> c;
  cin.ignore();
  cout << "Key is: " << c << endl;
  autokey cipher(c);
  cout << "Enter Plaintext : ";
  string text;
  getline(cin, text);
  string encrypted = cipher.encrypt(text);
  cout << "Encrypted text: " << encrypted << endl;
  string decrypted = cipher.decrypt(encrypted);
  cout << "Decrypted text: " << decrypted << endl;
  return 0;
}
\end{minted}

\section{Result and Discussion - I/O}
Using any other input other than an alphabet for the key loops back the user to enter a valid key, which should be an alphabet.
\begin{minted}{text} 
Enter a character key: 2
Key is: 2
Key is not an alphabet. Enter a new key: S
Enter Plaintext : One two three four five six seven eight.
Encrypted text: Gbr xpk hayvi jtil wndz waf pwzzr rmona.
Decrypted text: One two three four five six seven eight.
\end{minted}
Another example with a different key and plaintext which shows a problem with the encryption where the key is $A$ which is 0, and the string also contains $a$ or $A$ which is also 0, so the key does not change much in the string.
\begin{minted}{text}
Enter a character key: A
Key is: A
Enter Plaintext : abacadafagah                                
Encrypted text: abbccddffggh
Decrypted text: abacadafagah
\end{minted}

\section{Conclusion}
The Autokey cipher is more resistant to frequency analysis than the Vigen\`ere cipher because the plaintext itself generates the key stream. However, it remains vulnerable to key recovery attacks and is not used in modern cryptography.